\section{서론}

대화형 AI 에이전트를 위한 기존 벤치마크는 사용자와 효과적으로 소통하고 과제를 해결하기 위해 올바른 행동 순서를 수행하는 능력을 시험하도록 설계되어 있다~\cite{yao2024tau,lu2024toolsandbox, xiao2024flowbench, prabhakar2025apigen}. 이러한 벤치마크는 본질적으로 \emph{단일 제어(single-control)} 환경이다. AI 에이전트는 세계와 상호작용할 수 있지만, (시뮬레이션된) 사용자는 선호와 목표에 관한 정보를 제공하는 데 제한된다. 예를 들어 \taubench~\cite{yao2024tau}에서 \retail 및 \airline 도메인은 제약 충족(constraint satisfaction) 과제를 해결하는 에이전트의 능력을 시험한다. 여기서 제약은 에이전트가 따라야 하는 도메인 정책과, 믿음, 목표, 선호를 포함한 사용자의 인지 상태가 결합되어 발생한다.
\begin{figure}[tb]
    \centering
    \includegraphics[width=0.99\textwidth,trim=340 330 220 300,clip]{figures/raw/teaser.pdf}
    \vspace{-0.5em}
    \caption{\bench에서 지원하는 dual-control 환경. 에이전트는 데이터베이스와 상호작용하는 도구 집합에 접근할 수 있으며, 도메인 정책을 준수하면서 Tool-Agent-User (TAU) 상호작용을 통해 사용자의 요청을 해결해야 한다. 실제 시나리오를 시험하기 위해, 사용자는 시나리오 기반 지시와 자체 데이터베이스와 상호작용하는 도구 집합을 받은 또 다른 AI 에이전트로 시뮬레이션된다. 시뮬레이션 사용자는 TAU 상호작용을 이중 형식(Tool-User-Agent)으로 처리하는 더 쉬운 버전을 수행한다고 볼 수 있는데, 이때 사용자는 지시만 따르면 되며 과제 해결책을 추론할 필요는 없다.}
    \label{fig:teaser}
\end{figure}




이러한 설정에서 에이전트의 환경(에이전트가 취할 수 있는 행동 집합 포함)에 대한 사용자의 이해는 제약 충족 문제에 대해 하나의 해결 가능한 경로를 보장하도록 돕는 세심하게 작성된 자연어 지시문을 통해 제공된다. 이는 과제 명세를 더 쉽게 만들지만, 사용자는 의사소통을 통해서만 에이전트의 행동을 인지할 수 있고 초기 지시문에만 근거하여 환경 상태를 추론할 수 있다. 이는 기술 고객 지원과 같은 실제 시나리오와 상당히 다르다. 그런 경우 사용자는 요청을 받으면 문제를 진단하고 해결하기 위해 \textit{휴대전화를 재시작하기} 또는 \textit{비행기 모드 끄기} 같은 일부 행동을 능동적으로 수행해야 한다.

이러한 추가적인 실제 복잡성을 포착하기 위해 우리는 \textit{이중 제어(dual-control)} 환경을 갖춘 \bench를 소개한다. 여기서는 (LLM으로 시뮬레이션된) 사용자가 에이전트와 의사소통하는 것에 더해 행동을 취하고 도구를 호출할 수 있다. 이중 제어 모델은 선택적 정보 은닉, 환경의 비언어적 조작, 사용자의 과제나 시나리오를 더 쉽게 명세하는 것 등 사용자 시뮬레이션을 개선하는 여러 장점을 제공한다. 특히 도구의 구조화된 형식이 비구조화된 자연어보다 선호되는 경우에 그렇다. 

이중 제어 설정의 주된 과제는 도구를 통해 사용자 시뮬레이터에 의미 있는 행위 주체성(agency)을 부여하면서도 에이전트와 사용자 시뮬레이터 사이의 "복잡성 비대칭성"을 유지하는 데 있다. 즉 사용자의 어포던스(affordance)가 제한되어 있고, 사용자가 문제를 해결하려면 여전히 에이전트의 지원을 필요로 하도록 보장해야 한다. 사용자 시뮬레이터용 도구는 그 능력을 높여 이러한 균형을 깨뜨릴 수 있지만, 우리는 도구를 신중하게 설계하면 사용자 행동을 제약하여 더 신뢰할 수 있는 시뮬레이션으로 이어질 수 있음을 발견했다. 우리는 사용자 도구가 사람이 읽을 수 있는 출력만 산출하도록 보장하고, 사용자 계획을 에이전트 요청에 기반한 반응적 도구 사용으로 제한하며, 환경을 통해 사용자 행동을 긴밀하게 제약함으로써 이를 달성한다. 그 결과, 생각하는 동시에 사용자가 올바른 행동을 취하도록 효과적으로 안내해야 하는 에이전트를 위한 통제된 테스트베드가 만들어진다. 

전반적으로 \bench는 네 가지 핵심 기여를 한다:

\textbf{Telecom 이중 제어 도메인.} 우리는 \bench 안에 새로운 이중 제어 \telecom 도메인을 도입한다. 여기서는 이전 벤치마크와 달리 AI 에이전트와 사용자 모두 공유되고 동적인 환경의 상태를 관찰하고, 그 위에서 행동하며, 검증하기 위한 서로 다른 도구를 가진다. 이는 분산 부분 관찰 가능 마르코프 의사결정 과정(Decentralized Partially Observable Markov Decision Process, \decpomdp)~\cite{oliehoek2016concise}을 사용해 정식화된다. 이 이중 제어 설정은 실제 협업 시나리오를 정확히 나타내도록 설계되었으며, 사용자에게 제한된 행위 주체성만 있는 평가에는 없는 중요한 에이전트 조정 및 의사소통 과제를 드러낸다. 실험은 최첨단 LLM도 이 도메인에서 상당히 고전함을 보여 준다(예: 새 과제에서 \gptmodelshort \cite{gpt41}는 \passone 34\%, \ofourmodelshort \cite{o4mini}는 42\%, \claudesonnetshort \cite{claude37sonnet}는 49\%).

\textbf{조합적 과제 생성기.} \bench는 작은 원자적 기본 시나리오 집합(초기화, 해결, 검증 함수로 정의됨)으로부터 방대하고 다양한 검증 가능한 과제 집합을 자동으로 조합하는 프로그램적 과제 생성기를 통합한다. 이 방법은 과제의 증명 가능한 정확성을 보장하고, 완전한 도메인 포괄성을 제공하며, 과제 복잡도(예: 해결 단계 수 또는 문제 유형)에 대한 명시적 제어를 가능하게 하고, 손으로 만든 과제 모음과 관련된 수작업과 잠재적 취약성을 제거한다.

\textbf{신뢰할 수 있는 사용자 시뮬레이터.} 우리는 사용자 시뮬레이터를 환경과 긴밀히 결합하여 사용자 시뮬레이션의 신뢰성을 높인다. 사용자 행동은 이용 가능한 도구와 관찰 가능한 환경 상태에 의해 제약되며, 더 예측 가능하고 일관된 상호작용으로 이어진다. 이 접근은 사용자 시뮬레이터를 안내하기 위한 복잡한 자연어 프롬프팅의 필요성을 크게 줄이고, 훨씬 더 높은 신뢰성을 낳는다(예: \telecom 도메인의 사용자 시뮬레이터는 오류율 16\%, 치명적 오류 6\%를 보였는데, 이는 \taubench의 \retail 도메인에서 오류율 40\%, 치명적 오류 12\%였던 것과 대비된다).

 \textbf{에이전트 실패의 세분화된 진단.} \bench는 서로 다른 모드에서 과제 성공을 평가함으로써 에이전트 성능을 분해하여 진단할 수 있게 한다. (i) 에이전트가 모든 도구를 제어하여 추론 능력을 분리하는 완전 자율 모드(``no-user mode''), (ii) 의사소통 및 조정 요구사항을 도입하는 표준 이중 제어 모드가 그것이다. 우리의 결과는 에이전트가 자율 동작에서 사용자를 안내하는 방식으로 전환해야 할 때 성능이 크게 감소함(약 20\% \passone)을 드러낸다. 이는 순수 추론 실패와 의사소통 및 분산 제어에서 비롯된 실패를 명확히 구분하며, 후자가 현재 에이전트의 중요한 병목임을 정확히 지목한다.


